% Preamble bersama buku ajar (dimuat Quarto untuk PDF lewat include-in-header).
% Salin ke <proyek>/_bersama/preamble.tex dan sesuaikan warna seri.
% Mesin: LuaLaTeX, PDF bertag (\DocumentMetadata di _bersama/doc-class.tex).
% Tipografi: skill buku-ajar, references/aturan-gaya-dan-tampilan.md.
\usepackage{tikz}
\usetikzlibrary{bayesnet,calc,arrows.meta,positioning,shapes.geometric}
\usepackage{xcolor}
\definecolor{seriutama}{HTML}{0F3D3E}
\definecolor{seriabu}{HTML}{6B7280}

% Satu keluarga huruf: Libertinus Serif untuk teks dan judul, Libertinus Mono
% untuk kode, Libertinus Math untuk rumus (diatur di _quarto.yml). Judul KOMA
% memakai huruf teks tebal, bukan huruf sans.
\setkomafont{disposition}{\normalfont\bfseries\color{seriutama}}
\addtokomafont{chapter}{\Huge}
% Glif yang tidak ada di Libertinus Mono (mis. huruf Yunani di kode) diambil
% dari Libertinus Serif (fallback didefinisikan di _bersama/doc-class.tex).
\ifdefined\directlua\setmonofont{Libertinus Mono}[Scale=0.8,RawFeature={fallback=serimono}]\fi

% Blok lebar: kode, keluaran, dan gambar boleh lebih lebar daripada kolom
% prosa (kolom prosa dipersempit agar baris teks 40–70 karakter), tetapi
% tetap di dalam batas 18 mm dari tepi kertas. Di dalam kotak atau daftar
% (\linewidth < \textwidth) blok tidak dilebarkan.
\newlength{\tambahlebar}
\newlength{\tambahini}
\AtBeginDocument{\setlength{\tambahlebar}{\dimexpr(\paperwidth-36mm-\textwidth)/2\relax}}
% Tabel panjang belum didukung tagging LaTeX 2023: isinya tidak ditag agar
% struktur dokumen tetap sah (lihat KEPUTUSAN-RANCANGAN.md proyek).
\providecommand{\seritagstop}{\ifdefined\tagstop\tagstop\fi}
\providecommand{\seritagstart}{\ifdefined\tagstart\tagstart\fi}
\newcommand{\aturtambah}{%
  \ifdim\linewidth<\textwidth\setlength{\tambahini}{0pt}\fvset{breaklines,breakanywhere}%
  \else\setlength{\tambahini}{\tambahlebar}\fi}

% Ukuran huruf blok kode dan keluaran: Libertinus Mono 7,8 pt (Scale 0,8 di
% _quarto.yml dikalikan 9,75 pt) agar 72 karakter kode dan 78 karakter
% keluaran muat tanpa dipenggal.
\newcommand{\kodekecil}{\fontsize{9.75pt}{11.6pt}\selectfont}
% Di blok lebar, kode (maks. 72 karakter) dan keluaran (maks. 78) muat utuh.
% Di blok sempit (kotak, daftar) baris panjang dipenggal agar tidak melewati
% tepi; breaklines hanya dipakai di sana karena memecah tag paragraf PDF.
\usepackage{fvextra}
\DefineVerbatimEnvironment{Highlighting}{Verbatim}{commandchars=\\\{\},fontsize=\kodekecil}
\makeatletter
\AtBeginDocument{%
  % Keluaran R (verbatim) memakai blok lebar yang sama. Didefinisikan di sini
  % karena tcolorbox memuat paket verbatim yang menimpa definisi di preamble.
  \renewenvironment{verbatim}{\VerbatimEnvironment\aturtambah
    \begin{Verbatim}[fontsize=\kodekecil,
      xleftmargin=-\tambahini,xrightmargin=-\tambahini]}%
    {\end{Verbatim}}%
  % Kotak kode Quarto dibungkus blok lebar
  \ifdefined\Shaded
    \renewenvironment{Shaded}{\aturtambah\begin{tcolorbox}[borderline west={3pt}{0pt}{shadecolor},
      boxrule=0pt, breakable, colback={codebgcolor}, enhanced, frame hidden, sharp corners,
      left=2mm, right=1mm, top=1mm, bottom=1mm, boxsep=0pt,
      grow to left by=\tambahini, grow to right by=\tambahini]}{\end{tcolorbox}}%
  \fi
  % Gambar diperkecil hanya bila lebih lebar daripada blok lebar
  \ifdefined\pandocbounded
    \renewcommand*\pandocbounded[1]{%
      \aturtambah
      \sbox\pandoc@box{#1}%
      \Gscale@div\@tempa{\textheight}{\dimexpr\ht\pandoc@box+\dp\pandoc@box\relax}%
      \Gscale@div\@tempb{\dimexpr\linewidth+2\tambahini\relax}{\wd\pandoc@box}%
      \ifdim\@tempb\p@<\@tempa\p@\let\@tempa\@tempb\fi
      \makebox[\linewidth][c]{\ifdim\@tempa\p@<\p@\scalebox{\@tempa}{\usebox\pandoc@box}%
        \else\usebox{\pandoc@box}\fi}}%
  \fi
  % Tabel (longtable Pandoc) memakai blok lebar; lebar kolom relatif
  % dihitung dari \linewidth blok lebar
  \let\seriLTawal\longtable \let\endseriLTawal\endlongtable
  \renewenvironment{longtable}{\aturtambah
    \LTleft=-\dimexpr\tambahini\relax plus 1fill\relax
    \LTright=\dimexpr\tambahini\relax plus 1fill\relax
    \advance\linewidth 2\tambahini \advance\hsize 2\tambahini
    \par\seritagstop\seriLTawal}{\endseriLTawal\seritagstart}%
}
% Keterangan longtable Pandoc berada di \endhead sehingga kotak kepala tabel
% (beserta \write dari \label dan anchor tautannya) disalin ulang di setiap
% halaman lanjutan: label terdefinisi ganda (telaah pakar seri Bayesian,
% 2026-11-07; dicegat G-LOG). Kepala asli menjadi kepala halaman pertama
% (\LT@firsthead); salinan untuk halaman lanjutan dibersihkan dari whatsit
% write dan pdf_dest (LuaTeX). Keterangan tetap diulang di halaman lanjutan.
\ifdefined\directlua
\directlua{
  function seri_buang_tulis(n)
    local WH, HL, VL = node.id("whatsit"), node.id("hlist"), node.id("vlist")
    local hapus = {[node.subtype("write")] = true, [node.subtype("pdf_dest")] = true}
    local function bersih(head)
      local cur = head
      while cur do
        local nxt = cur.next
        if cur.id == WH and hapus[cur.subtype] then
          head = node.remove(head, cur); node.free(cur)
        elseif cur.id == HL or cur.id == VL then
          cur.list = bersih(cur.list)
        end
        cur = nxt
      end
      return head
    end
    local b = tex.getbox(n)
    if b then b.list = bersih(b.list) end
  end
}
% \endhead = \LT@end@hd@ft\LT@head (longtable 4.17) dengan penyalinan kepala
% sebelum \LT@bchunk membuka potongan tabel berikutnya.
\AtBeginDocument{%
  \def\endhead{%
    \LT@echunk
    \ifx\LT@start\endgraf
      \LT@err{Longtable head or foot not at start of table}{Increase LTchunksize}%
    \fi
    \setbox\LT@head\box\z@
    \ifvoid\LT@firsthead
      \setbox\LT@firsthead\copy\LT@head
      \directlua{seri_buang_tulis(\number\LT@head)}%
    \fi
    \LT@get@widths
    \LT@bchunk}}
\fi
\makeatother

% Baris yatim dan janda dicegah (Permendikbudristek 22/2022, Lampiran III:
% bebas baris tunggal; telaah pakar seri Bayesian 2026-11-07)
\widowpenalty=10000
\clubpenalty=10000
\displaywidowpenalty=10000

% Cegah tabel pendek dan blok kode terbelah halaman (\needspace di bab)
\usepackage{needspace}
% Pemenggalan EYD V dan nama yang tidak dipenggal (dibangkitkan harness/penggalan.py)
\IfFileExists{../_bersama/penggalan-id.tex}{\input{../_bersama/penggalan-id.tex}}{}
% Sampul (opsional): halaman sampul dimuat sebelum halaman judul Quarto.
% Sampul dikompilasi dengan LuaLaTeX dan huruf Libertinus Serif/Math.
% \usepackage{pdfpages}
% \providecommand{\halamansampul}{1}
% \AtBeginDocument{\includepdf[pages=\halamansampul]{../sampul/sampul-seri.pdf}}
% Tanda air draf di setiap halaman (keputusan penulis 2026-11-07).
% Hapus dua baris ini untuk versi terbit.
\usepackage{draftwatermark}
\DraftwatermarkOptions{text={DRAF}, fontsize=120pt, angle=45, color={gray!14}}
% Paragraf menjorok, tanpa spasi antarparagraf, termasuk paragraf pertama
% setelah judul (keputusan penulis 2026-11-07; pasangannya indent: true di _quarto.yml)
\usepackage{indentfirst}
\setlength{\parindent}{1.5em}
\setlength{\parskip}{0pt plus 1pt}

% Penyesuaian tagging (latex-lab, TL2023) untuk judul KOMA-Script: judul
% subbab menjadi struktur H2/H3 yang tertutup dan paragraf sesudahnya P.
\usepackage{etoolbox}
\makeatletter
\ExplSyntaxOn
\cs_if_exist:NT \__tag_setup_restore_para_default:
  {
    \cs_new_protected:Npn \buku_judul_awal:n #1 { \tagtool{para-flattened=true} }
    \clist_map_inline:nn { section, subsection, subsubsection, paragraph, subparagraph }
      {
        \AddtoDoHook{heading/begingroup/#1}{\buku_judul_awal:n}
        \AddtoDoHook{heading/branch/star/#1}{\buku_judul_awal:n}
      }
    \pretocmd{\@xsect}
      { \tag_tool:n {restore-para} \bool_set_false:N \l__tag_para_flattened_bool }
      {}{ \PackageWarning{bukuajar}{\string\@xsect\space tidak dapat ditambal} }
    \renewcommand{\sectionlinesformat}[4]
      {
        \IfUseNumber
          { \@kernel@tag@hangfrom{\hskip #2#3}{#4} }
          { \@hangfrom{\hskip #2#3}{#4} }
      }
  }
\ExplSyntaxOff
% Judul bab yang memuat koma dipecah latex-lab menjadi beberapa kunci
% (sec-start-chapter=#2, lalu title=#1); argumen dibungkus tiga lapis kurung
% kurawal (satu dilepas xparse, satu per tingkat keyval).
\ifcsname tag_tool:n\endcsname
  \NewCommandCopy\seribabasli\chapter
  \RenewDocumentCommand\chapter{s o m}{%
    \IfBooleanTF{#1}%
      {\IfNoValueTF{#2}{\seribabasli*[{{{#3}}}]{#3}}{\seribabasli*[{{{#2}}}]{#3}}}%
      {\IfNoValueTF{#2}{\seribabasli[{{{#3}}}]{#3}}{\seribabasli[{{{#2}}}]{#3}}}}
\fi
% Kode sebaris (\texttt) boleh dipatahkan sesudah _ / . < > = [ ( , tanpa
% tanda hubung, agar nama fungsi dan berkas panjang tidak melewati tepi kolom.
% Dikerjakan di Lua (LuaLaTeX): penalti 200 disisipkan sesudah karakter itu
% bila hurufnya Libertinus Mono dan karakter berikutnya masih huruf kode.
\ifdefined\directlua
\directlua{
  local patah = {[0x5F]=true, [0x2F]=true, [0x2E]=true, [0x3C]=true, [0x3E]=true,
                 [0x3D]=true, [0x5B]=true, [0x28]=true, [0x2C]=true}
  local GLYPH = node.id("glyph")
  local mono = {}
  local function ismono(f)
    if mono[f] == nil then
      local d = font.getfont(f)
      mono[f] = d ~= nil and d.name ~= nil and string.find(d.name, "LibertinusMono", 1, true) ~= nil
    end
    return mono[f]
  end
  local function sisip(head)
    for n in node.traverse_id(GLYPH, head) do
      local s = n.next
      if patah[n.char] and ismono(n.font) and s and s.id == GLYPH and ismono(s.font) then
        local p = node.new("penalty")
        p.penalty = 200
        node.insert_after(head, n, p)
      end
    end
    return head
  end
  luatexbase.add_to_callback("pre_linebreak_filter", sisip, "seri.kodepatah")
}
\fi
% Kolom prosa lebih sempit: beri kelonggaran spasi lebih besar sebelum baris luap
\setlength{\emergencystretch}{5em}
% Nomor halaman tiga digit muat di daftar isi
\renewcommand*{\@pnumwidth}{2.4em}
% Rumus tampilan memakai blok lebar yang sama dengan kode (di luar kotak).
\newcommand{\serirumuslebar}{%
  \ifdim\hsize=\textwidth
    \displayindent=\dimexpr\displayindent-\tambahlebar\relax
    \displaywidth=\dimexpr\displaywidth+2\tambahlebar\relax
  \fi}
\AtBeginDocument{\everydisplay\expandafter{\the\everydisplay\serirumuslebar}}
\makeatother
